% Préambule commun aux cours IPR Scholar.
\documentclass[11pt,a4paper,oneside]{article}
\usepackage[T1]{fontenc}
\usepackage[utf8]{inputenc}
\usepackage[french]{babel}
\usepackage{lmodern}
\usepackage{microtype}
\usepackage[a4paper,margin=2.35cm]{geometry}
\usepackage{amsmath,amssymb,mathtools}
\usepackage{booktabs,longtable,array}
\usepackage{xcolor}
\usepackage{hyperref}
\usepackage{tikz}
\usetikzlibrary{arrows.meta,positioning,shapes.geometric,calc}
\hypersetup{colorlinks=true,linkcolor=blue!55!black,citecolor=blue!55!black,urlcolor=blue!55!black}
\definecolor{iprblue}{HTML}{1F4E79}
\definecolor{iprgreen}{HTML}{2D6A4F}
\newcommand{\IPR}{\textsc{ipr}}
\newcommand{\coursmeta}[2]{\author{Projet \IPR{} -- documentation Scholar}\date{#1}\maketitle\begin{center}\small #2\end{center}\vspace{0.5cm}\tableofcontents\clearpage}
\newcommand{\resultat}[1]{\paragraph{À retenir.} #1}

\setcounter{tocdepth}{1}
\hypersetup{pdftitle={Cours 10 -- Filtre de Haar et ondelettes},
pdfauthor={Projet IPR -- Scholar}}
\title{Cours 10 -- Filtre de Haar et ondelettes\\
De la somme et de la différence à l'estimation du bruit photographique}
\begin{document}
\coursmeta{5 octobre 2026}{Cours théorique, explicatif et applicatif ; exemples,
exercices corrigés et liens avec le moteur actuel de \IPR{}.}

\section{Objectifs et parcours de lecture}
Ce cours part d'une paire de nombres, construit une transformation réversible,
puis l'étend à une image. L'objectif final est de comprendre pourquoi
\texttt{(p00-p01-p10+p11)*0.5} et \texttt{0.67448975} apparaissent dans
le moteur, et ce qu'ils permettent réellement de mesurer.
Les prérequis sont les sommes, les matrices élémentaires, la moyenne et la variance.
Lire d'abord les exemples ; les démonstrations peuvent être reprises ensuite.

\resultat{Haar n'est ni une IA ni, à lui seul, un débruiteur.
C'est un changement de représentation : il sépare des informations de
fond et des variations locales. Un traitement des coefficients est nécessaire
pour modifier effectivement le signal.}

\section{Intuition : conserver le fond, isoler la variation}
Deux pixels voisins $u=10$ et $v=14$ contiennent une luminosité moyenne
de 12 et une différence de 4. En conservant ces deux informations, on peut
retrouver les deux pixels. Leur moyenne seule perdrait de l'information.
Pour comparer des énergies et des variances sans correction supplémentaire,
Haar emploie une normalisation orthonormale :
\begin{equation}\label{eq:pair}
a=\frac{u+v}{\sqrt{2}},\qquad d=\frac{u-v}{\sqrt{2}}.
\end{equation}
$a$ est l'\emph{approximation} et $d$ le \emph{détail}. Le signe de $d$
indique quel pixel est le plus grand, dans la convention choisie ici.
$a$ n'est pas la moyenne arithmétique : cette dernière vaut $a/\sqrt{2}$.
Pour notre paire : $a=12\sqrt{2}$, $d=-2\sqrt{2}$.

\begin{center}
\begin{tikzpicture}[node distance=1.2cm,box/.style={draw=iprblue,rounded corners,
align=center,minimum height=1cm}]
\node[box] (pair) {Deux pixels\\$(u,v)$};
\node[box,right=of pair] (haar) {Somme et différence\\normalisées};
\node[box,above right=0.15cm and 1cm of haar] (approx) {Approximation $a$\\fond local};
\node[box,below right=0.15cm and 1cm of haar] (detail) {Détail $d$\\variation locale};
\draw[-{Latex}] (pair)--(haar);
\draw[-{Latex}] (haar)--(approx);
\draw[-{Latex}] (haar)--(detail);
\end{tikzpicture}
\end{center}

\section{Construction mathématique en une dimension}
\subsection{Filtres passe-bas et passe-haut}
Dans notre convention de produit scalaire sur une paire, les filtres sont
\begin{equation}
\ell=\frac{1}{\sqrt{2}}(1,1),\qquad
h=\frac{1}{\sqrt{2}}(1,-1).
\end{equation}
Le premier additionne et privilégie les variations lentes ; le second
soustrait et annule une constante. Appliquer ces filtres puis garder un
coefficient par paire correspond à un filtrage suivi d'une décimation par deux.
Une bibliothèque utilisant une convolution peut inverser le filtre ou changer
le signe des détails ; l'inverse doit suivre la même convention.

Pour un signal $x[n]$ de longueur paire :
\begin{equation}
a[k]=\frac{x[2k]+x[2k+1]}{\sqrt{2}},\qquad
d[k]=\frac{x[2k]-x[2k+1]}{\sqrt{2}}.
\end{equation}
Chaque sortie a deux fois moins d'échantillons, mais les deux sorties réunies
conservent le nombre initial de valeurs. Il n'y a pas de compression sans
une étape supplémentaire de quantification ou de suppression de coefficients.

\subsection{Orthogonalité, inverse et énergie}
Écrivons $\binom{a}{d}=Q\binom{u}{v}$, avec
\begin{equation}
Q=\frac{1}{\sqrt{2}}\begin{pmatrix}1&1\\1&-1\end{pmatrix}.
\end{equation}
Les lignes ont une norme unité et un produit scalaire nul.
Ainsi $Q^{\mathsf T}Q=I$ et l'inverse est
\begin{equation}
u=\frac{a+d}{\sqrt{2}},\qquad v=\frac{a-d}{\sqrt{2}}.
\end{equation}
En développant les carrés, on obtient $a^2+d^2=u^2+v^2$.
Cette conservation concerne l'énergie quadratique, pas l'énergie physique
des photons. La lecture de Haar comme une base orthogonale est également
présentée par Strang~\cite{strangHaar}. Les calculs et exemples de ce cours
sont des développements pédagogiques propres au projet.

\subsection{Lecture fréquentielle}
Avant décimation, les réponses sont
\begin{equation}
L(\omega)=\frac{1+e^{-i\omega}}{\sqrt{2}},\quad
H(\omega)=\frac{1-e^{-i\omega}}{\sqrt{2}}.
\end{equation}
Leurs modules valent respectivement
$\sqrt{2}|\cos(\omega/2)|$ et $\sqrt{2}|\sin(\omega/2)|$.
Le passe-haut est nul à fréquence zéro ; le passe-bas est nul à $\pi$.
Ces filtres très courts n'ont pas la séparation fréquentielle d'un filtre
idéal. La reconstruction exacte provient de la banque complète, pas d'un
passe-bas parfait ni de la conservation d'une seule branche.

\section{Pourquoi parler d'ondelette et de plusieurs échelles ?}
Une ondelette est localisée : elle regarde une zone limitée du signal.
Pour Haar, la fonction d'échelle $\phi$ vaut 1 sur $[0,1)$ et 0 ailleurs.
L'ondelette mère est
\begin{equation}
\psi(t)=\begin{cases}
1,&0\le t<1/2,\\-1,&1/2\le t<1,\\0,&\text{ailleurs}.
\end{cases}
\end{equation}
Elle a une intégrale nulle : une constante ne produit aucun détail.
Les fonctions $\psi_{j,k}(t)=2^{j/2}\psi(2^j t-k)$ changent d'échelle et
de position. La transformée discrète implémente cette idée avec les paires.
On redécompose l'approximation, pas nécessairement toutes les branches.

\subsection{Exemple complet sur quatre valeurs}
Pour $x=(4,6,10,12)$ :
\begin{align}
a_1&=(5\sqrt{2},11\sqrt{2}),&d_1&=(-\sqrt{2},-\sqrt{2}),\\
a_2&=16,&d_2&=-6.
\end{align}
La représentation finale est $(a_2,d_2,d_1)$ : quatre valeurs au total.
En remontant l'inverse à deux niveaux, on retrouve exactement $x$ en
arithmétique exacte. Vérification énergétique :
\begin{equation}
4^2+6^2+10^2+12^2=296=16^2+(-6)^2+2+2.
\end{equation}
Pour une longueur $N=2^J$, la somme du travail sur les niveaux est
$N+N/2+N/4+\cdots<2N$ : complexité linéaire pour cette construction.
Les dimensions non dyadiques nécessitent une politique de bord ou un
nombre de niveaux limité ; ce n'est pas une raison de redimensionner la photo.

\section{Extension à une image : quatre informations par bloc}
Prenons $P=\begin{pmatrix}p_{00}&p_{01}\\p_{10}&p_{11}\end{pmatrix}$,
où le premier indice désigne la ligne et le second la colonne.
Appliquer Haar sur les lignes puis sur les colonnes donne :
\begin{align}
A&=\frac{p_{00}+p_{01}+p_{10}+p_{11}}{2},\\
D_x&=\frac{p_{00}-p_{01}+p_{10}-p_{11}}{2},\\
D_y&=\frac{p_{00}+p_{01}-p_{10}-p_{11}}{2},\\
D_{xy}=HH&=\frac{p_{00}-p_{01}-p_{10}+p_{11}}{2}.\label{eq:hh}
\end{align}
$A$ est deux fois la moyenne des quatre pixels ; $D_x$ mesure les différences
gauche/droite, $D_y$ les différences haut/bas, $HH$ une différence croisée.
La notion de détail \emph{diagonal} ne signifie pas détection de tous les
contours diagonaux, ni exclusion des autres structures.

Les noyaux correspondants rendent les signes explicites :
\begin{equation}
\frac12\begin{pmatrix}+1&+1\\+1&+1\end{pmatrix},\quad
\frac12\begin{pmatrix}+1&-1\\+1&-1\end{pmatrix},\quad
\frac12\begin{pmatrix}+1&+1\\-1&-1\end{pmatrix},\quad
\frac12\begin{pmatrix}+1&-1\\-1&+1\end{pmatrix}.
\end{equation}
Les noms LH/HL et horizontal/vertical dépendent des conventions d'axes.
La documentation PyWavelets précise les siennes et fournit
\texttt{dwt2}, \texttt{idwt2} et la décomposition multi-niveaux~\cite{pywtHaar}.
Pour relire le code, préférer les équations aux seuls noms de sous-bandes.

\subsection{Forme matricielle et reconstruction}
Pour $p=(p_{00},p_{01},p_{10},p_{11})^{\mathsf T}$ :
\begin{equation}
\begin{pmatrix}A\\D_x\\D_y\\HH\end{pmatrix}
=T p,\qquad
T=\frac12\begin{pmatrix}
1&1&1&1\\1&-1&1&-1\\1&1&-1&-1\\1&-1&-1&1
\end{pmatrix}.
\end{equation}
$T^{\mathsf T}T=I$ ; on reconstruit par $p=T^{\mathsf T}(A,D_x,D_y,HH)^{\mathsf T}$.
Par exemple $p_{00}=(A+D_x+D_y+HH)/2$.
En 2D, chaque sous-bande a un quart des pixels d'entrée : les quatre
sous-bandes réunies conservent la taille totale. Plusieurs niveaux
redécomposent la sous-bande d'approximation.

\subsection{Exemple numérique vérifiable}
Pour $P=\begin{pmatrix}10&12\\14&18\end{pmatrix}$,
$(A,D_x,D_y,HH)=(27,-3,-5,1)$.
Le premier pixel reconstruit est $(27-3-5+1)/2=10$.
L'énergie vaut $764$ aussi bien dans les pixels que dans les coefficients :
$10^2+12^2+14^2+18^2=27^2+(-3)^2+(-5)^2+1^2$.

\section{Du traitement du signal à la physique du bruit}
\subsection{Pourquoi HH annule certains fonds}
Une plage constante donne $HH=0$. Pour un plan affine
$p(x,y)=c+\alpha x+\beta y$, la même annulation se produit :
les différences croisées retirent le fond et les deux pentes.
Cela favorise l'analyse du bruit sur des aplats ou des dégradés réguliers.
Ce n'est pas une annulation de toute image lisse : courbures mixtes,
textures et frontières peuvent rester dans HH.

\subsection{Démonstration du facteur 1/2}
Supposons quatre bruits indépendants, centrés, de même variance $\sigma^2$.
Le bruit du coefficient HH est leur combinaison de poids
$w=(1,-1,-1,1)/2$. Sa variance vaut
\begin{equation}
\operatorname{Var}(HH)=\sum_i w_i^2\sigma^2
=4\times\frac14\sigma^2=\sigma^2.
\end{equation}
Le facteur 1/2 maintient donc l'échelle du bruit. Avec 1/4, la variance
serait $\sigma^2/4$ et l'écart-type serait divisé par deux ; une correction
supplémentaire serait nécessaire pour estimer le bruit d'entrée.

\subsection{Quand le modèle simple cesse d'être exact}
Si les variances diffèrent, mais que les pixels restent indépendants :
\begin{equation}
\operatorname{Var}(HH)=\frac14\sum_{i=1}^4\sigma_i^2.
\end{equation}
Avec une covariance $\Sigma$ quelconque :
\begin{equation}
\operatorname{Var}(HH)=w^{\mathsf T}\Sigma w.
\end{equation}
Un bruit identique ajouté aux quatre pixels s'annule complètement :
HH peut être nul alors que chaque pixel est bruité.
Le filtre mesure donc une composante spatiale du bruit, pas tout le bruit.

Le bruit photonique dépend du signal : le nombre de photons suit un modèle
de Poisson idéal, dont moyenne et variance sont égales. Un modèle simple
après conversion peut s'écrire $\operatorname{Var}(n\mid s)\simeq\alpha s+\beta$.
Les paramètres dépendent du capteur et de la chaîne de lecture. Dématriçage,
corrections optiques et filtrage introduisent ensuite des corrélations.
Un unique MAD global sur une image développée ne retrouve pas à lui seul
ces paramètres physiques. Haar ne dépend pas directement de l'ISO.

\section{MAD : passer des coefficients à un estimateur robuste}
Pour des coefficients $z_i=HH_i$, définir
\begin{equation}
m=\operatorname{median}(z_i),\qquad
\operatorname{MAD}=\operatorname{median}(|z_i-m|).
\end{equation}
La médiane est moins sensible aux valeurs isolées extrêmes que la moyenne ;
elle n'est pas immunisée contre une texture couvrant une grande partie des blocs.
Le MAD n'est ni la moyenne des distances ni le MAD des intensités brutes.
SciPy documente cette statistique et sa normalisation gaussienne~\cite{scipyMad}.

Pour $Z\sim\mathcal N(0,1)$ et $q>0$ :
\begin{equation}
\Pr(|Z|\le q)=2\Phi(q)-1.
\end{equation}
La médiane de $|Z|$ impose $2\Phi(q)-1=1/2$, donc
$q=\Phi^{-1}(3/4)\simeq0{,}67448975$.
D'où l'estimateur, cohérent dans ce modèle à grand échantillon :
\begin{equation}\label{eq:mad}
\widehat\sigma=\frac{\operatorname{MAD}}{0{,}67448975}.
\end{equation}
La constante est statistique, pas une calibration esthétique.
Un petit échantillon peut donner une estimation imprécise ; hors modèle
gaussien indépendant, le résultat reste surtout un proxy de dispersion.

\subsection{Exemple et interprétation d'un ratio}
Pour $z=(-2,-1,0,1,2)$, $m=0$, les distances sont $(2,1,0,1,2)$,
le MAD vaut 1 et $\widehat\sigma\simeq1{,}4826$.
Si le proxy passe de $0{,}010$ à $0{,}006$, le ratio vaut 0,6.
Sous un modèle inchangé, cela correspond à 40\,\% de baisse d'écart-type
et 64\,\% de baisse de variance, puisque $0{,}6^2=0{,}36$.
Ce n'est pas une mesure de 40\,\% de pixels nettoyés ni une preuve de qualité.

\section{Estimer n'est pas débruiter : le seuillage}
Une transformée suivie de son inverse, sans modification, restitue le signal.
Un débruiteur à ondelettes transforme, réduit des détails puis reconstruit.
Le seuillage doux est
\begin{equation}
S_\lambda(d)=\operatorname{sgn}(d)\max(|d|-\lambda,0).
\end{equation}
Par exemple $S_2(-5)=-3$ et $S_2(1)=0$.
Le seuillage dur conserve $d$ si $|d|>\lambda$ et donne zéro sinon.
Il est discontinu au seuil ; le seuillage doux est continu mais réduit
aussi l'amplitude des coefficients conservés.

Les travaux de Donoho étudient cette famille dans un modèle gaussien
et sous des hypothèses précises~\cite{donohoSoft}. Pour une transformée
orthonormale de vecteur discret, un seuil universel classique est
$\lambda=\widehat\sigma\sqrt{2\log M}$, avec $M$ le nombre de coefficients
de bruit considérés. La normalisation de la base doit rester cohérente.
Ce seuil peut trop lisser une photographie : ce n'est pas le seuil de \IPR{}.
Avec $\widehat\sigma=0{,}01$ et $M=65536$, il vaut environ $0{,}0471$.

\section{Application exacte au projet IPR}
\subsection{Moteur : \texttt{estimateDenoiseSigma}}
Dans \texttt{src/core/DetailProcessing.cpp}, le moteur :
\begin{enumerate}
\item construit les plans Y Rec.709 linéaire, $Co=R-B$ et $Cg=G-(R+B)/2$ ;
\item prélève des blocs \emph{adjacents} $2\times2$, avec un pas dans la grille
de blocs pour viser un budget de l'ordre de 16384 échantillons ;
\item exclut les blocs exactement constants, pour éviter qu'une bordure noire
impose un MAD nul au reste de l'image ;
\item calcule HH, sa médiane, le MAD centré et la division par 0,67448975.
\end{enumerate}
Les dimensions impaires laissent la dernière ligne/colonne hors de cette
estimation. Un plan trop petit ou sans bloc retenu donne une estimation nulle.
Le budget 16384 est un compromis de coût, pas une constante physique.
Exclure les blocs constants conditionne aussi la population mesurée :
ce choix ne garantit pas un estimateur sans biais du bruit de toute l'image.
La médiane C++ utilise le rang supérieur lorsque l'effectif est pair ;
NumPy utilise la moyenne des deux rangs centraux. Les estimations peuvent
donc différer légèrement, même avant les autres différences de sélection.

\subsection{Analyse des TIFF : \texttt{noise\_sigma}}
Dans \texttt{tests/analyse\_noise.py}, les régions contiguës de 256 pixels
sont lues en 16 bits, décodées sRGB vers le linéaire, puis analysées.
Le script utilise tous les blocs complets de la région, y compris les blocs
constants. Il applique le même HH et le MAD centré, mais n'est pas une copie
bit-exacte de l'estimateur C++ global. Il faut distinguer le diagnostic sur
une région exportée du paramètre estimé avant les étapes finales du pipeline.

\subsection{Ce que le moteur ne fait pas}
Le débruitage courant n'effectue \emph{pas} une décomposition/reconstruction
Haar complète. Haar fournit le niveau estimé du bruit ; le lissage réel
travaille sur deux résidus gaussiens, avec guide de structure et seuillage.
Le seuil fin est $\lambda=\widehat\sigma m s$, avec $s$ le curseur divisé
par 100, $m=2{,}55$ pour Y ou $3{,}15$ pour la couleur.
Ces multiplicateurs sont empiriques. Ils ne doivent pas être confondus avec
les constantes normalisées 1/2 ou 0,67448975.
Les corrections de moyenne, de gamut et le mélange progressif sont des
étapes supplémentaires, pas des propriétés automatiques de Haar.

\begin{center}
\begin{tabular}{lll}
\toprule Constante & Nature & Rôle \\ \midrule
$1/2$ & Normalisation & Norme unité du noyau HH \\
$0{,}67448975$ & Statistique & Quantile normal pour calibrer le MAD \\
16384 & Calcul & Budget approximatif d'échantillonnage \\
$10^{-7}$ & Numérique & Estimation trop petite : débruitage ignoré \\
2,55 / 3,15 & Empirique & Intensité maximale du seuil fin \\
\bottomrule
\end{tabular}
\end{center}
Le seuil numérique $10^{-7}$ est dans les unités du plan flottant concerné ;
ce n'est pas un niveau de bruit universel du capteur.

\section{Limites à reconnaître dans les photographies}
\begin{longtable}{p{0.25\linewidth}p{0.65\linewidth}}
\toprule Situation & Conséquence \\ \midrule
Aplat idéal + bruit blanc & Le MAD HH peut bien estimer l'écart-type.\\
Cheveux, pelage, feuillage & De vraies variations entrent dans les coefficients ; risque de sur-estimation.\\
Bruit spatialement corrélé & La variance HH n'est plus nécessairement celle des pixels.\\
Grandes taches colorées & HH peut peu réagir ; une petite statistique n'implique pas absence de bruit couleur.\\
Quantification ou zones noires & Beaucoup de coefficients identiques peuvent donner un MAD nul ou des plateaux.\\
Pas régulier de prélèvement & Certaines structures périodiques peuvent être sur- ou sous-représentées.\\
Décalage d'un pixel & La grille des paires change ; transformée décimée non invariante par translation.\\
Dimensions impaires & L'estimation locale actuelle omet un bord ; une transformée complète nécessite une règle explicite.\\
\bottomrule
\end{longtable}

Une frontière verticale idéale peut produire $D_x\ne0$ mais $HH=0$.
Inversement un damier peut produire de forts HH sans bruit aléatoire.
Si tous les HH du damier sont égaux, le MAD \emph{centré} vaut pourtant zéro.
Le MAD décrit la dispersion des coefficients, pas leur amplitude maximale.

Haar est court, peu coûteux et facilement explicable ; ses fonctions en
marches peuvent donner des artefacts en blocs si on les utilise pour un
débruitage complet agressif. Des filtres plus longs ou une transformation
non décimée peuvent être envisagés, au prix d'autres compromis. Les valeurs
de cette section sont des conséquences des équations, pas un classement
expérimental d'algorithmes pour le corpus IPR.

\section{Travaux pratiques reproductibles}
Le fichier \texttt{references/haar-exemples.py} accompagne ce cours.
Il ne lit ni ne modifie aucune photographie et demande seulement NumPy :
\begin{verbatim}
python docs/scholar/references/haar-exemples.py
\end{verbatim}
Il vérifie les exemples numériques, l'inverse et l'énergie, puis génère
un aplat gaussien déterministe pour comparer $\sigma$ connu au MAD HH.
Il montre aussi qu'un bruit commun aux quatre pixels peut disparaître de HH.
Les tirages synthétiques ne sont pas une validation de qualité photographique.
Convertir toujours les entiers TIFF en flottants \emph{avant} soustraction :
en entier non signé, un résultat négatif peut déborder.

L'extrait de principe, pour un plan flottant de dimensions paires, est :
\begin{verbatim}
hh = (p[::2, ::2] - p[::2, 1::2]
      - p[1::2, ::2] + p[1::2, 1::2]) / 2
center = np.median(hh)
sigma = np.median(np.abs(hh - center)) / 0.67448975
\end{verbatim}
Le script accompagnant ajoute les validations d'entrée et les assertions ;
ce fragment seul n'est pas une API de production.

\section{Exercices corrigés}
\subsection*{1. Une paire constante}
Calculer Haar pour $(7,7)$ puis reconstruire.
\paragraph{Correction.} $a=7\sqrt{2}$, $d=0$ ; l'inverse redonne $(7,7)$.
Un détail nul ne signifie pas un signal nul.

\subsection*{2. Un bord vertical sans bruit}
Pour $P=\begin{pmatrix}0&20\\0&20\end{pmatrix}$, calculer les quatre coefficients.
\paragraph{Correction.} $(A,D_x,D_y,HH)=(20,-20,0,0)$.
La frontière est présente malgré un coefficient diagonal nul.

\subsection*{3. Pourquoi ne pas diviser HH par quatre ?}
Quelle serait la variance sur un bruit indépendant de variance $\sigma^2$ ?
\paragraph{Correction.} Les quatre poids seraient $\pm1/4$ ; la somme
de leurs carrés vaut $1/4$. L'écart-type estimé serait donc divisé par deux
si l'on gardait la calibration MAD actuelle sans correction.

\subsection*{4. MAD et valeur extrême}
Calculer le MAD de $(-1,0,0,1,100)$.
\paragraph{Correction.} Médiane 0 ; distances $(1,0,0,1,100)$ ; MAD 1.
La valeur 100 influence peu cette statistique, sans prouver que le signal
est gaussien ni que la valeur extrême est un pixel chaud.

\subsection*{5. Lecture critique d'un rapport}
Un proxy passe de 0,02 à 0,01 après filtrage. Peut-on conclure que
la moitié des détails ont été conservés et que le rendu est deux fois meilleur ?
\paragraph{Correction.} Non. Le ratio d'écart-type proxy est 0,5 ;
le ratio de variance associé serait 0,25 sous les mêmes hypothèses.
Il faut contrôler luminosité, textures, couleur et contours séparément.

\section{Glossaire et ressources du projet}
\begin{description}
\item[Filtre] Combinaison pondérée de valeurs voisines.
\item[Ondelette] Fonction localisée, déplacée et changée d'échelle pour analyser un signal.
\item[Approximation] Composante passe-bas ; pas nécessairement une moyenne non normalisée.
\item[Détail] Composante de variation, pouvant contenir structure et bruit.
\item[Décimation] Conservation d'un échantillon sur un nombre fixé après filtrage.
\item[Orthogonalité] Produit scalaire nul entre deux directions de représentation.
\item[Orthonormalité] Orthogonalité et norme unité ; conservation de l'énergie quadratique.
\item[HH] Détail obtenu par le filtre passe-haut dans les deux directions.
\item[MAD] Médiane des distances absolues à la médiane.
\item[Proxy] Indicateur indirect, dont les hypothèses et limites doivent rester explicites.
\end{description}
Ce cours complète \path{cours-08-bruit-grain-signal.tex}.
Les ressources, depuis la racine du dépôt, sont :
\begin{itemize}
\item contrat : \path{docs/architecture/denoising-sharpening.md} ;
\item moteur : \path{src/core/DetailProcessing.cpp} ;
\item analyseur : \path{tests/analyse_noise.py} ;
\item campagne : \path{docs/validation/analyses/analyse-bruit-v2-2026-10-05.md}.
\end{itemize}
Le cours explique des modèles ; les preuves de validation restent dans
les rapports et données brutes, et le moteur n'est pas modifié par ce document.

\resultat{Dans IPR, Haar révèle des variations fines et le MAD en estime
la dispersion. Leurs constantes normalisées ont une justification mathématique,
mais l'efficacité du débruitage photographique dépend aussi du modèle de bruit,
de la sélection spatiale, des seuils, des textures et de la reconstruction.}

\bibliographystyle{plain}
\begingroup\raggedright\bibliography{references-haar}\endgroup
\end{document}
